\documentclass[11pt]{article}

\usepackage[T1]{fontenc}
\usepackage[utf8]{inputenc}
\usepackage{lmodern}
\usepackage[a4paper,margin=28mm]{geometry}
\usepackage{microtype}
\usepackage{amsmath,amssymb,mathtools}
\usepackage{enumitem}
\usepackage{booktabs,tabularx}
\usepackage[colorlinks=true,linkcolor=blue,citecolor=blue,urlcolor=blue]{hyperref}

\title{A Projection-First Reframing of Information and Computation:\\
\large Distinguishability, Channels, and Conditional Undecidability}
\author{Peter Nero}
\date{July 2026\\Version 2}

\begin{document}
\maketitle

\begin{abstract}
A projection can identify source states that an effective description no
longer distinguishes. This observation is useful, but it is not by itself a
theory of information, computation, complexity, or undecidability.
Information quantities require a state or probability law and a specified
channel; computation requires an encoding, evolution rule, readout, and
resource model; undecidability requires an unbounded problem family and an
explicit reduction from a known undecidable problem. This paper gives a
projection-first account of these layers. Classical and quantum
data-processing inequalities formalize loss of distinguishability under
channels, while recovery maps state when the loss can be reversed on a
selected family of states. Noninvertibility alone does not imply entropy
production, an arrow of time, computational hardness, or undecidability.
For Modal Triplet Theory, finite selected operators supply concrete
descriptive channels, but a physical undecidability theorem remains open
until a uniform computational embedding is constructed.
\end{abstract}

\section*{Version 2 Revision Note}
\begin{description}[leftmargin=3.2cm,style=nextline]
\item[Supersedes] Version 1.0, \emph{A Projection-First Reframing of
Information, Computation, and Undecidability}.
\item[Reason] The original paper inferred information loss, complexity, and
generic undecidability directly from projection and finite admissibility.
Those implications do not hold without further probabilistic, dynamical,
and computational structures.
\item[Resolution] Projection is now typed as a channel only after a state
space is supplied. Information loss is formulated through data processing
and recovery. Computation and undecidability are made conditional on
explicit encodings, resource models, and reductions.
\item[Retained result] Projection-first language remains useful for
separating source distinctions from distinctions available to an effective
observer.
\item[Remaining boundary] MTT still needs a selected uniform computational
embedding, together with a proved reduction, before it can claim a new
physical undecidability result.
\end{description}

\tableofcontents

\section{Scope}

Let
\[
 \pi:X\longrightarrow Y
\]
map a richer source description to an effective one. If
\(\pi(x)=\pi(x')\), the effective variable does not distinguish \(x\) from
\(x'\). This is a statement about an equivalence relation induced by
\(\pi\). It is not yet a statement about entropy, information flow, memory,
or computation.

Those notions belong to different mathematical categories. A probability
distribution is required before Shannon entropy is defined. A density
operator and observable algebra are required before quantum information is
defined. An input language and an algorithmic model are required before
computability or complexity is defined. The purpose of this paper is to
retain the useful projection-first intuition while making those additional
structures explicit.

\section{From a map to an information channel}

\subsection{Classical distinguishability}

Suppose \(X\) is a random variable with distribution \(p\), and let
\(Y=\pi(X)\). The deterministic map induces a stochastic channel
\[
 K(y\mid x)=
 \begin{cases}
 1,&y=\pi(x),\\
 0,&\text{otherwise}.
 \end{cases}
\]
Only now are Shannon entropy and mutual information defined. For a
deterministic channel on finite alphabets,
\[
 H(Y)\leq H(X).
\]
The inequality describes the chosen random variables and their law. It is
not a property of the bare set map independent of \(p\).

A more robust statement compares two candidate source laws \(p\) and \(q\).
For any stochastic channel \(K\), the classical data-processing inequality
gives
\[
 D(Kp\Vert Kq)\leq D(p\Vert q),
\]
where \(D\) is relative entropy. Operationally, processing cannot improve
the ability to distinguish \(p\) from \(q\). A projection therefore supports
an information-loss interpretation only relative to specified states and
an observation channel.

\subsection{Quantum distinguishability}

In quantum theory the corresponding object is a completely positive,
trace-preserving map
\[
 \mathcal E:\mathcal T(\mathcal H_X)\longrightarrow
 \mathcal T(\mathcal H_Y).
\]
Quantum relative entropy obeys
\[
 D(\mathcal E(\rho)\Vert\mathcal E(\sigma))
 \leq D(\rho\Vert\sigma).
\]
A finite projector \(P\) does not by itself define this channel on all
states. One must specify, for example, whether the operation is a selected
outcome, the nonselective instrument
\[
 \rho\longmapsto P\rho P+(I-P)\rho(I-P),
\]
or a compression followed by normalization. These operations have
different physical meanings.

\section{Loss, recovery, and irreversibility}

Information being unavailable in \(Y\) does not imply that it was destroyed
in \(X\). A section \(s:Y\to X\) satisfying \(\pi s=I_Y\) merely chooses one
representative from each fiber. It does not recover the original \(x\) from
\(\pi(x)\).

For a channel \(\mathcal E\), recovery on a selected state family
\(\mathcal S\) requires another channel \(\mathcal R\) such that
\[
 \mathcal R\mathcal E(\rho)=\rho
 \qquad(\rho\in\mathcal S),
\]
or an explicit norm or fidelity error bound. Equality in suitable
data-processing inequalities is tied to such sufficient or recoverable
families; approximate equality motivates quantitative recovery estimates
\cite{Petz1986,Sutter2016}.

This distinction prevents three common overclaims.
\begin{enumerate}[leftmargin=*]
\item A many-to-one description can coexist with reversible lower dynamics.
\item Entropy need not increase along a trajectory unless a state,
coarse-graining rule, and dynamical law make it do so.
\item Memory effects require a history-dependent reduced law, hidden
variables, or a non-Markovian channel; they do not follow from
noninjectivity alone.
\end{enumerate}

Physical irreversibility therefore needs a semigroup, dissipative
generator, boundary condition, growing recovery error, or another
time-directed structure. Projection can explain which distinctions are
discarded, but not by itself why their loss is dynamically irreversible.

\section{What makes a physical process a computation?}

A computational interpretation needs at least a tuple
\[
 (\mathcal I,E,\Phi_t,R,\mathcal O),
\]
where \(\mathcal I\) is an input language, \(E\) encodes inputs into physical
states, \(\Phi_t\) evolves those states, \(R\) reads an output, and
\(\mathcal O\) specifies the output convention. Complexity additionally
requires an input-size function, an error tolerance, and a resource cost
such as time, memory, energy, or circuit depth.

Projection may enter this tuple in several useful ways. It can define a
readout, compress inaccessible variables, or identify physically equivalent
encodings. It can also make a chosen decoding problem ill posed. None of
these facts alone establishes that the underlying task is algorithmically
hard. A constant map is maximally noninjective but trivial to compute.
Conversely, an injective map may be expensive to evaluate.

The phrase ``the universe computes'' is therefore optional interpretation,
not a mathematical consequence. The precise question is whether a physical
family implements a named input-output problem with controlled errors and
resources.

\section{Undecidability requires a reduction}

Undecidability is stronger than practical unpredictability, chaos, or an
expensive numerical calculation. A decision problem is undecidable when no
algorithm halts with the correct answer for every input in a specified
unbounded family.

A physical undecidability proof consequently needs:
\begin{enumerate}[leftmargin=*]
\item a recursively describable family of physical instances \(M_w\);
\item a decision predicate \(Q(M_w)\);
\item an effective map from instances \(w\) of a known undecidable problem;
\item a proof that \(Q(M_w)\) answers that problem; and
\item robustness conditions showing that the encoding belongs to the
claimed physical class.
\end{enumerate}
This is the pattern used in genuine many-body undecidability results: a
Hamiltonian family is constructed so that its spectral-gap behavior encodes
the halting problem \cite{Cubitt2015,Bausch2020}.

By contrast, the following do not prove undecidability:
\begin{itemize}
\item a projection being many-to-one;
\item a finite stability margin;
\item sensitivity to initial data;
\item a simulation taking a long time;
\item not knowing in advance whether a numerical solver will converge; or
\item one fixed finite matrix having a complicated spectrum.
\end{itemize}
A fixed finite exact model is, in principle, exhaustively decidable for
finite predicates. Undecidability can enter only through a uniform
unbounded family, an infinite-volume limit, an exact real-number oracle, or
another explicitly stated source of unbounded computation.

\section{Consequences for Modal Triplet Theory}

MTT currently supplies several concrete finite descriptive structures:
selected projectors, finite Hessian blocks, a \(27\times27\) finite carrier
at its declared Standard-Model profile tier, and selected free-field
operator data. These can be studied as channels once states and instruments
are specified. Their finite nature is an advantage for reproducibility, but
it does not support a new undecidability claim.

A serious MTT computational program has two distinct branches.

\paragraph{Finite branch.}
Specify the state, selected instrument, readout, and error metric for each
finite operator. Then compute distinguishability loss, recovery fidelity,
spectral conditioning, and algorithmic cost. This can produce exact or
certified numerical results without invoking undecidability.

\paragraph{Uniform branch.}
Construct a family indexed by words, lattice size, cutoff, bundle data, or
another recursive parameter. Prove that the family preserves the selected
MTT constraints and embeds a universal computation. Only then should a
halting, reachability, spectral, or admissibility predicate be tested for an
undecidability reduction.

The second branch is open. It may succeed, but projection and admissibility
do not replace the construction.

\section{Claim ledger}

\begin{center}
\small
\begin{tabularx}{\textwidth}{@{}p{0.30\textwidth}p{0.18\textwidth}X@{}}
\toprule
Claim & Status & Required structure or qualification\\
\midrule
\(\pi\) identifies source states & Exact & A defined map \(X\to Y\).\\
Projection reduces distinguishability & Conditional & State family and
classical or quantum channel.\\
Data processing gives monotonicity & Standard result & Relative entropy and
an admissible stochastic or CPTP map.\\
Projection destroys information & Not implied & Must exclude source access
and recovery on the selected family.\\
Projection creates irreversibility & False in general & Needs directed
dynamics, boundary data, or recovery-error growth.\\
Projection creates computational hardness & False in general & Needs a
problem encoding and resource model.\\
Finite MTT operators are undecidable & Not established & A fixed finite
instance is not an undecidable family.\\
MTT admits physical undecidability & Open & Uniform selected family and an
explicit reduction.\\
\bottomrule
\end{tabularx}
\end{center}

\section{Discussion}

The corrected projection-first view is narrower than the original one, but
more useful. Projection says which distinctions survive a change of
description. Information theory quantifies those distinctions after a
state and channel are supplied. Recovery theory tests whether the loss is
reversible on a chosen family. Computation theory asks whether a controlled
encoding and readout implement a problem. Undecidability appears only when
an unbounded family carries a proved reduction.

This hierarchy also clarifies black-hole and measurement language.
Tracing out a region, conditioning on an outcome, recording an apparatus
state, and losing access to a source are different channels. Calling all of
them ``projection'' can hide the very mechanism under study. MTT should
therefore name the channel and recovery criterion in each application.

\section{Conclusion}

Projection-first reasoning does not make information, computation, or
undecidability inevitable. It gives a disciplined starting point: identify
the source states, the effective states, and the distinctions that the map
forgets. The next layers must then be built explicitly.

For MTT, the immediate rigorous opportunity is finite and quantitative:
turn selected projectors and operators into specified channels and compute
their distinguishability and recovery properties. A new undecidability
result would be a later theorem, requiring a recursive physical family and
a genuine computational reduction.

\begin{thebibliography}{9}

\bibitem{Shannon1948}
C.~E. Shannon,
\emph{A Mathematical Theory of Communication},
Bell System Technical Journal \textbf{27} (1948), 379--423, 623--656.

\bibitem{Turing1936}
A.~M. Turing,
\emph{On Computable Numbers, with an Application to the
Entscheidungsproblem},
Proceedings of the London Mathematical Society \textbf{42} (1936),
230--265.

\bibitem{Petz1986}
D. Petz,
\emph{Sufficient Subalgebras and the Relative Entropy of States of a von
Neumann Algebra},
Communications in Mathematical Physics \textbf{105} (1986), 123--131.

\bibitem{Sutter2016}
D. Sutter, M. Tomamichel, and A.~W. Harrow,
\emph{Strengthened Monotonicity of Relative Entropy via Pinched Petz
Recovery Map},
IEEE Transactions on Information Theory \textbf{62} (2016), 2907--2913.

\bibitem{Cubitt2015}
T.~S. Cubitt, D. P\'erez-Garc\'ia, and M.~M. Wolf,
\emph{Undecidability of the Spectral Gap},
Nature \textbf{528} (2015), 207--211.

\bibitem{Bausch2020}
J. Bausch, T.~S. Cubitt, A. Lucia, and D. P\'erez-Garc\'ia,
\emph{Undecidability of the Spectral Gap in One Dimension},
Physical Review X \textbf{10} (2020), 031038.

\end{thebibliography}

\end{document}
